\documentclass[11pt,a4paper]{article}
\usepackage[utf8]{inputenc}
\usepackage[T1]{fontenc}
\usepackage[french]{babel}
\usepackage{lmodern}
\usepackage[margin=2cm]{geometry}
\usepackage{microtype}
\usepackage{amsmath}
\usepackage{booktabs,tabularx,array}
\usepackage{xcolor}
\usepackage{enumitem}
\usepackage{hyperref}
\definecolor{rfmgreen}{HTML}{276B50}
\hypersetup{colorlinks=true,urlcolor=rfmgreen,linkcolor=rfmgreen,pdftitle={Segmentation clients RFM},pdfauthor={Groupe Togo AI Summer School}}
\setlength{\parindent}{0pt}
\setlength{\parskip}{5pt}
\setlist[itemize]{label=\textbullet,nosep,leftmargin=*}
\setlist[enumerate]{nosep,leftmargin=*}
\newcolumntype{Y}{>{\raggedright\arraybackslash}X}
\newcommand{\fichier}[1]{\texttt{\small\detokenize{#1}}}

\begin{document}
\begin{center}
{\LARGE\bfseries Segmentation clients RFM}\par
{\large Mini-rapport synthétique}\par
{\color{rfmgreen}Togo AI Summer School}\par
\end{center}

\textbf{Membres du groupe :} AKAYA Eniwinéwé Henri ; ALKISSANKDEI T. Djamal ;
BLAISE Mouné Tchoubou ; LARE V. Donné ; METO Kemi Gédéon ;
OPEKOU-DOUDOE Christian ; TORA Dkawlma Emmanuel.

\textbf{Origine du projet.} L'analyse initiale résulte d'un travail collectif
réalisé lors du Togo AI Summer School. À partir de ce socle, le porteur du dépôt
a poursuivi le projet à titre personnel : approfondissement des diagnostics,
ajout d'une API FastAPI, d'un assistant n8n avec mémoire PostgreSQL et d'une
interface web. Ces extensions sont distinguées du travail initial du groupe.

\section{Objectif et méthode}
L'objectif est de regrouper les clients selon leurs achats observés afin de
formuler des hypothèses de campagnes différenciées. La segmentation décrit
l'historique ; elle ne prédit ni le départ d'un client ni sa rentabilité.

Les deux feuilles du jeu \textit{UCI Online Retail II} sont concaténées. Le
nettoyage supprime les doublons exacts, les lignes sans identifiant client ou
date exploitable, les factures d'annulation et les quantités ou prix non
positifs du périmètre principal. Les retours sont conservés séparément pour
l'analyse de sensibilité. Les achats retenus couvrent le
\textbf{1\textsuperscript{er} décembre 2009 au 9 décembre 2011}, sur
\textbf{5\,878 clients}. La référence est le \textbf{10 décembre 2011}.

Pour chaque client $i$, les indicateurs sont :
\begin{itemize}
\item $R_i$ : nombre de jours depuis son dernier achat valide à la date de référence ;
\item $F_i$ : nombre de factures positives distinctes dans la fenêtre ;
\item $M_i$ : somme des montants d'achats positifs, en livres sterling (GBP), sans déduction des retours.
\end{itemize}
Une faible récence signifie un achat récent, pas nécessairement un nouveau
client. Le montant ne mesure pas la marge, car les coûts ne sont pas disponibles.

Les trois variables sont plafonnées à leur 99\textsuperscript{e} percentile,
transformées par $\log(1+x)$ puis standardisées :
\[
 z_{ij}=\frac{\log\!\left(1+\min(x_{ij},q_{0,99,j})\right)-\mu_j}{\sigma_j}.
\]
Les paramètres sont appris sur la population historique. Cette préparation
réduit l'influence des valeurs extrêmes et évite que les montants dominent
mécaniquement les distances. K-means utilise une graine fixée à 42 et
100 initialisations. La livraison documente scikit-learn 1.9.0.

\section{Choix du nombre de segments}
Les valeurs de $k$ de 2 à 8 sont comparées avec l'inertie, la silhouette,
Calinski-Harabasz, Davies-Bouldin, les tailles des groupes et la stabilité
ARI entre dix graines. La silhouette, Calinski-Harabasz et Davies-Bouldin
préfèrent ici \textbf{$k=2$}. Leur compromis par moyenne des rangs préfère
également 2 ; le coude géométrique indicatif suggère \textbf{$k=4$}.

Le choix final de \textbf{$k=5$ est commercial}, et non un optimum statistique
démontré. Les profils et la matrice de correspondance entre $k=4$ et $k=5$
permettent d'examiner les différences observées. L'intérêt d'une campagne
supplémentaire reste à mesurer face à une stratégie à quatre groupes.

\newpage
\section{Résultats et propositions marketing}
Le montant total des achats positifs retenus est de
\textbf{17\,685\,460,64 GBP}. Les parts ci-dessous utilisent ce même dénominateur.
Les noms des groupes sont relatifs à cette partition et décrivent leurs profils
moyens ; ils ne caractérisent pas uniformément chaque client du groupe.

\begin{center}
\small
\renewcommand{\arraystretch}{1.35}
\begin{tabularx}{\textwidth}{@{}Yrrrr@{}}
\toprule
\textbf{Segment} & \textbf{Clients} & \textbf{CA (\%)} & \textbf{R (j)} & \textbf{F} \\
\midrule
Champions & 587 & 59,18 & 18,00 & 28,56 \\
Actifs à montant élevé & 1\,250 & 24,12 & 58,36 & 8,56 \\
Achats anciens & 1\,335 & 10,10 & 303,15 & 3,56 \\
Achats récents occasionnels & 1\,126 & 4,35 & 32,30 & 2,50 \\
Achats anciens à faible montant & 1\,580 & 2,25 & 416,98 & 1,23 \\
\bottomrule
\end{tabularx}
\end{center}
{\small R et F : moyennes de récence et de fréquence. Les pourcentages sont arrondis.
Source : \fichier{tableau_synthese_segments.csv}.}

\begin{center}
\small
\renewcommand{\arraystretch}{1.25}
\begin{tabularx}{\textwidth}{@{}Yr@{}}
\toprule
\textbf{Segment} & \textbf{Montant moyen (GBP)} \\
\midrule
Champions & 17\,829,92 \\
Actifs à montant élevé & 3\,412,10 \\
Achats anciens & 1\,337,73 \\
Achats récents occasionnels & 683,75 \\
Achats anciens à faible montant & 252,15 \\
\bottomrule
\end{tabularx}
\end{center}

Les Champions représentent environ 10\,\% des clients et 59,18\,\% du montant
positif. Cette concentration justifie une attention particulière, sans prouver
une marge supérieure ou une fidélité future. Les achats anciens ne constituent
pas, à eux seuls, une preuve d'attrition. De même, les achats récents
occasionnels ne renseignent pas la date du premier achat du client.

\textbf{Actions proposées, à tester plutôt qu'à considérer comme acquises :}
\begin{enumerate}
\item \textbf{Champions :} tester un programme relationnel et des accès anticipés.
Mesurer le réachat et les désabonnements face à un groupe témoin.
\item \textbf{Actifs à montant élevé :} proposer des produits complémentaires.
Comparer le montant supplémentaire et le coût de contact ; une hausse de CA
ne suffit pas à démontrer une hausse de marge.
\item \textbf{Achats anciens :} tester une relance sur un échantillon et comparer
le taux de réachat à celui d'un groupe témoin.
\item \textbf{Achats récents occasionnels :} encourager un prochain achat par
une sélection pertinente, sans supposer qu'il s'agit du deuxième achat.
\item \textbf{Achats anciens à faible montant :} limiter la pression de contact
et tester une relance peu coûteuse. Mesurer le réachat, les désabonnements et
le coût par réponse avant généralisation.
\end{enumerate}

Pour comparer les stratégies à quatre et cinq segments, les groupes de test
devraient être randomisés, avec un budget, un horizon et une pression de
contact comparables. Le bénéfice commercial ne sera établi qu'après mesure,
par exemple de la conversion ou de la marge incrémentale si les coûts sont
collectés. Aucun résultat de campagne n'est disponible dans ce projet.

\newpage
\section{Vérification de l'effet des retours}
Une seconde segmentation utilise le montant net sur les \textbf{mêmes clients},
avec la même fenêtre, les mêmes R et F et $k=5$. Les lignes à quantité négative
et prix positif sont déduites du montant, qu'elles portent ou non un numéro
de facture commençant par C. Les annulations à quantité positive restent
ambiguës et ne sont pas transformées arbitrairement en montants négatifs.
Les clients ayant uniquement des retours sont hors de la population comparée.

Les plafonds et la standardisation appris sur les achats positifs restent
figés. La transformation du montant net conserve son signe :
\[
 g(M)=\operatorname{sign}(M)\log(1+|M|).
\]
La valeur absolue intervient uniquement dans le logarithme : les montants
nets négatifs ne deviennent pas des dépenses positives. K-means est
réentraîné sur cette variante. Les groupes sont appariés en maximisant les
clients communs avant de compter les migrations.

Les retours retenus représentent \textbf{1\,031\,016,71 GBP}, soit
\textbf{5,83\,\%} des achats positifs. L'ARI entre les deux partitions vaut
\textbf{0,5810} et \textbf{1\,671 clients (28,43\,\%)} changent de groupe
après appariement. Les retours modifient donc la partition dans ce protocole.
Les répétitions sur dix graines documentent la sensibilité aux initialisations ;
elles ne constituent ni un intervalle de confiance ni un test de significativité.

Cette expérience isole l'effet monétaire : les retours ne sont pas rapprochés
facture par facture et peuvent concerner des achats antérieurs à la fenêtre.
Des clients sans retour peuvent aussi migrer lorsque les centres changent.
Les noms alignés des groupes nets sont des repères, pas des profils métier validés.

\section{Extensions de la plateforme et limites}
L'API sert les résultats et classe un profil RFM par proximité aux centroïdes
figés, sans réentraînement ni ajout aux données. Les trois dates de la fenêtre
et de référence sont obligatoires. Le mode historique exige les dates du
modèle ; une autre période impose une simulation explicite, sans annualisation
ni probabilité de confiance. La cohérence des dates ne garantit pas la
pertinence sur une autre population.

L'interface permet la recherche, le filtrage et l'export global des clients.
L'assistant n8n consulte les résultats agrégés et utilise une mémoire
PostgreSQL associée à la session. Une recette de 14 scénarios couvre notamment
les chiffres, la concision, la mémoire, les pannes et les injections. Les tests
techniques ne remplacent pas une recette réelle des réponses du fournisseur.
La synchronisation contrôle le modèle, les CSV et la fiche documentaire avant livraison.

Sans vérité terrain, les indices internes, la stabilité et les profils ne
prouvent pas l'efficacité commerciale. La saisonnalité, l'ancienneté des
clients, les grossistes et les choix de nettoyage peuvent influencer les
résultats. L'usage doit limiter les sollicitations, permettre le retrait,
protéger les données individuelles et éviter les traitements discriminatoires
ou les prix personnalisés injustifiables.

\textbf{Traçabilité.} Résultats issus de \fichier{ml/index.ipynb}, des tables de
\fichier{ml/outputs/} et de \fichier{docs/current-model.md}.
Modèle : \fichier{rfm-3ab5423591563489}. Livraison :
\fichier{bundle-229a31b580c84328}. Ce rapport constitue une synthèse de cette
livraison ; ses chiffres doivent être revus après un nouvel entraînement.
\end{document}
